\begin{lemma}[A coefficient-independent recursive elimination bound]
\label{c:recursive-elimination-bound}
\status{Proved}
For integers $v,d\geq1$, let $\mathcal B(v,d)$ be the tree of
finite sequences $(\alpha_0,\ldots,\alpha_{r-1})$ in $\Z_{\geq0}^v$
satisfying
\[
 |\alpha_i|\leq d2^i,
 \qquad \alpha_j\not\leq\alpha_i\quad(0\leq j<i<r),
\]
where $\leq$ is coordinatewise.  This is a finite tree.  Its
maximum sequence length $\mathcal L(v,d)$ is a total computable
integer-valued function.  Put
\[
 \mathcal D(v,d)=d\,2^{\mathcal L(v+1,d)}.
\]
Over any field, an ideal in $v$ ordinary polynomial variables
generated in degrees at most $d$ has, for every monomial order,
a Gr\"obner basis in degrees at most $\mathcal D(v,d)$.
In particular this bound applies to an elimination order.
\end{lemma}

\begin{proof}
We give the elementary argument, so no closed Gr\"obner-degree
estimate is needed.  Dickson's assertion that an infinite
sequence in $\Z_{\geq0}^v$ contains a comparable pair follows
by induction on $v$.  An infinite sequence of nonnegative
integers has an infinite nondecreasing subsequence: either some
value repeats infinitely or one successively selects larger
values.  Apply this observation to the first coordinate and
then apply induction to the remaining coordinates.  Thus
$\mathcal B(v,d)$ has no infinite branch.  Each node has only
finitely many children, by its displayed degree bound.
If the tree were infinite, repeatedly choosing a child with
infinitely many descendants would construct an infinite branch.
It is consequently finite.  Enumerating its successive finite
levels until the next level is empty computes $\mathcal L(v,d)$.

Homogenize the given generators with one new variable $h$.
Extend the desired monomial order by comparing total degree
first and, at equal degree, comparing the monomials after
$h=1$ in the desired order.  Process the homogeneous generators
one at a time, reducing each by the previous nonzero remainders,
and then run the usual pair-reduction algorithm.  Every accepted
nonzero remainder has a leading monomial not divisible by any
previously accepted leading monomial.  Homogeneous division
does not increase degree, and the least common multiple of
two leading monomials has degree at most the sum of their
degrees.  If the accepted polynomials are numbered from zero,
their degrees are therefore at most $d2^i$.  Their leading
exponent vectors form a sequence in $\mathcal B(v+1,d)$.
At most $\mathcal L(v+1,d)$ polynomials can be accepted, and
all their degrees are at most $\mathcal D(v,d)$.
There are then only finitely many pairs to process.

For completeness, the pair criterion here says that if every
pair polynomial reduces to zero, every element of the generated
ideal has leading monomial divisible by one of the recorded
leading monomials.  Indeed write such an element as a sum of
monomial multiples of the recorded polynomials.  If the greatest
displayed leading monomial cancels, subtract the corresponding
pair relations to lower that greatest monomial.  A representation
minimizing it, and then the number of terms attaining it, cannot
have such a cancellation.  Its greatest monomial is therefore
the leading monomial of the element.  This proves the criterion
and that the resulting list is a Gr\"obner basis.

Dehomogenization gives a Gr\"obner basis for the original ideal.
To check this even though the homogenized generator ideal need
not be saturated, if $f$ belongs to the original ideal, then
$h^e f^{\mathrm h}$ belongs to the homogenized generator ideal
for some $e\geq0$, by homogenizing an expression for $f$ in
the generators to a common degree.  A homogeneous leading
monomial of our basis divides that of $h^e f^{\mathrm h}$.
After setting $h=1$, it divides the leading monomial of $f$
in the desired order.  Degrees do not increase.  This proves
the bound for arbitrary, including elimination, orders.
\end{proof}

\begin{corollary}[A uniform arity bound for rational cohomology tables]
\label{c:uniform-arity-bound}
\status{Proved in the stated rational-table scope}
Suppose the graded Weil cohomology algebra in
Theorem~\ref{c:finite-unconditional-criterion} has a rational
table.  This includes rational Betti cohomology over $\C$.
Put $q=2n$, $u=a-1$, $v=b$, $s=u+v$, and
\[
 M_0=\frac12\sum_{m=1}^{q}\bigl(s^m+(u-v)^m\bigr),\qquad
 M_1=\frac12\sum_{m=1}^{q}\bigl(s^m-(u-v)^m\bigr).
\]
Define the following integers using the function of the preceding
lemma:
\[
 O=M_1+4uv,\qquad
 E=M_0+2(u^2+v^2)+2,\qquad
 J=2^O,\qquad P=E+J,
\]
\[
 D=\mathcal D(P,q+2),\qquad
 B_{n,a,b}=\max\{2,\ q+1,\ 2a(b+1),\ q(D-1+M_1)\}.
\]
Then all the point-and-diagonal cohomological relations hold in
rational Chow on all powers if and only if the cycle map is
injective on $T_{B_{n,a,b}}^{\CH}(X,o)$.
The bound depends only on $(n,a,b)$, and is a total computable
integer.  It is not a practical or minimal bound, and does not
decide the required Chow injectivity.
\end{corollary}

\begin{proof}
We bound the input arities of the finite form relations in
Theorem~\ref{c:finite-unconditional-criterion}.  Its parameter
superspace has $M_0$ even and $M_1$ odd coordinates.
Present $\operatorname{GL}(u|v)$ by two independent even
supermatrices $G,H$ with relations $GH=HG=1$; $H$ is the inverse
of $G$.  This presentation has $2(u^2+v^2)$ even and $4uv$ odd
matrix entries.  Present the independent line torus by $z,w$
with $zw=1$.  The orbit-graph equation for a form of arity $m$
is a coordinate of the parameter tensor minus a coordinate of
$z(H^\vee)^{\otimes m}t_m$, with the fixed dual-action signs.
It has total degree at most $m+1\leq q+1$; all group equations
have degree two.  Because the table is rational, its structure
constants add no variables or degrees.

There are $E$ even and $O$ odd variables in this graph
presentation.  The free supercommutative algebra is a free module
of rank $J=2^O$ over the polynomial ring $A_0$ on its even
variables, with exterior-monomial basis $e_1,\ldots,e_J$.
If $f_1,\ldots,f_t$ are the graph equations, their ideal is the
$A_0$-submodule generated by all $e_j f_i$.  The polynomial
coordinates of these module generators have even-variable
degree at most $q+1$.  Encode a vector $(p_1,\ldots,p_J)$ by
the ordinary polynomial $\sum p_jY_j$, and form the ideal in
$A_0[Y_1,\ldots,Y_J]$ generated by the encoded module generators
and all $Y_iY_j$.  Its part of total $Y$-degree one is exactly
the original module.  This ordinary ideal has $P=E+J$ variables
and generators of degree at most $q+2$.

Choose an elimination order that eliminates the even group
variables and every $Y_j$ whose exterior monomial involves an
odd group variable.  Retain the even parameter variables and
the remaining $Y_j$.  The ordinary ideal is homogeneous in
total $Y$-degree.  The homogeneous pair-reduction construction in
Lemma~\ref{c:recursive-elimination-bound} preserves this additional
grading, both before and after dehomogenization, and therefore
gives a bounded Gr\"obner basis homogeneous in $Y$-degree.  The elimination basis elements of
$Y$-degree one give module generators for the intersection with
the original parameter superalgebra, namely the schematic
super-orbit ideal.  Quadratic $Y$ terms cannot contribute to
this degree-one part.

Lemma~\ref{c:recursive-elimination-bound} bounds these ordinary
basis elements by $D$.  Their coefficients in the retained even
variables have degree at most $D-1$.  Returning to exterior
monomials adds at most $M_1$ odd parameter factors.  Every
parameter coordinate has input arity at most $q$.  The orbit
ideal therefore has generators of input arity at most
$q(D-1+M_1)$.  Taking homogeneous components and closing their
span under the supergroup preserves input arity.  The split
pure tensor modules used to define $P_\nu$ can consequently
all be taken within this same bound.  The assertion now follows
from Corollary~\ref{c:finite-power-cutoff}.
\end{proof}

\begin{remark}[The coefficient qualification in the uniform bound]
\label{c:uniform-arity-scope}
\status{Proved scope clarification}
The finite criterion and finite cut-off of
Theorem~\ref{c:finite-unconditional-criterion} and
Corollary~\ref{c:finite-power-cutoff} apply to an arbitrary
Weil field.  The displayed bound depending only on $(n,a,b)$
has the additional rational-table hypothesis.  For a presented
coefficient algebra whose $h$ named generators include every
coordinate of $t=(\kappa_1,\ldots,\kappa_q)$, with its exact
defining ideal generated in degrees at most $e$, the same proof gives a bound by adding
$h$ to $E$ and replacing the generator-degree input $q+2$
by $\max(q+3,e+1)$.  The extra one accounts for a structure
constant occurring as a variable in an orbit equation.  These
$h$ coefficient variables are eliminated together with the group
variables, so the result is the rational contraction $I_\Q$.
We make no uniform rational-contraction claim for unspecified
elements of an arbitrary coefficient field.

The accompanying exact certificate exhausts only
$\mathcal L(1,d)=d+1$ for $1\leq d\leq4$ and
$\mathcal L(2,1)=4$, and checks the square-zero elimination
mechanism for the explicit super map
$z=x$, $\theta_1=\eta$, $\theta_2=x\eta$.
It does not evaluate the enormous geometric cut-off above.
No closed Dub\'e degree formula is used in its proof.
\end{remark}
